\begin{scanpage}{0092}{75}
\[
\left.\begin{aligned}
j_x^0&=-\frac c{2\pi}E_{0x}\sin\gamma\sin\psi\,e^{ik\Phi},\\
j_y^0&=\frac c{2\pi}E_{0x}\sin\gamma\cos\psi\,e^{ik\Phi},\\
j_z^0&=\frac c{2\pi}E_{0x}\cos\gamma\cos\psi\,e^{ik\Phi},
\end{aligned}\right\}\tag{13.02}
\]
and with the H-polarization it has the components
\[\left.\begin{aligned}j_x^0=j_y^0&=0,\\j_z^0&=-\frac c{2\pi}H_{0x}\sin\psi\,e^{ik\Phi},\end{aligned}\right\}\tag{13.03}\]
where
\[\Phi=a\sin\gamma\sin\psi+\zeta\cos\gamma.\tag{13.04}\]
Let us calculate the field created by these currents in the region $\varphi=-\frac\pi2$.

The vector potential of the fringing field is determined by the equations
\[\mathbf A=\frac ac\int_0^{2\pi}d\psi\int_{-l/2}^{l/2}\mathbf j^0(\zeta,\psi)\frac{e^{ikr}}r\,d\zeta\quad\text{with }\gamma=0\tag{13.05}\]
and
\[\mathbf A=\frac ac\int_{-\pi}^0d\psi\int_{-l/2}^{l/2}\mathbf j^0(\zeta,\psi)\frac{e^{ikr}}r\,d\zeta\quad\text{with }\gamma>0,\tag{13.06}\]
where
\[r=\sqrt{\xi^2+(y-\eta)^2+(z-\zeta)^2}.\tag{13.07}\]
Since the field in the far zone $(R\gg ka^2,\ R\gg kl^2)$ is of interest to us, these expressions may be simplified by using the relationship
\[r\simeq R+a\sin\vartheta\sin\psi-\zeta\cos\vartheta.\tag{13.08}\]
\end{scanpage}
